\section{\ref{sec:sleep}장. 수면}

수면 모드, 재생, 시냅스 축소는 뉴런 기록, 미세투석, 전자 현미경으로 측정되었다. 수면 일정과 느린 시소는 이 책의 모델로 기술했고, 수면의 목적은 대부분 가설이다.

{\footnotesize
\setlength{\tabcolsep}{3pt}\renewcommand{\arraystretch}{1.15}
\begin{longtable}{>{\raggedright}p{0.5cm}>{\raggedright}p{4.0cm}>{\raggedright}p{5.6cm}>{\raggedright}p{2.3cm}>{\raggedright\arraybackslash}p{1.7cm}}
\toprule
No. & 명제 & 실험과 측정한 것 & 출처 & 상태 \\
\midrule\endfirsthead
\toprule
No. & 명제 & 실험과 측정한 것 & 출처 & 상태 \\
\midrule\endhead
1 & 잠자는 동안 피질 뉴런은 조용하지 않다. 바뀌는 것은 동작 모드다 & 쥐 피질 뉴런의 장시간 기록: 서파 수면에서 발화율은 0이 아니며, 활동 기간과 침묵 기간이 번갈아 온다. 오래 깨어 있은 뒤에는 모든 상태에서 발화율이 높고, 잔 뒤에는 낮다 & \cite{Vyazovskiy2009} & 측정됨 \\
2 & \nt{ACET}은 낮과 렘수면에서 높고 서파 수면에서 낮다(표~\ref{tab:sleepmodes}) & 자유롭게 움직이는 고양이의 피질 미세투석: 각성 시 아세틸콜린 방출이 서파 수면보다 $100$--$175\%$ 높고, 렘수면에서는 조용한 각성과 거의 같은 수준이다 & \cite{Marrosu1995} & 측정됨 \\
3 & \nt{NORA}와 \nt{SERO}는 서파 수면에서 잠잠해지고 렘수면에서 거의 조용하다 & 쥐의 개별 \nt{NORA} 뉴런과 고양이의 \nt{SERO} 뉴런을 수면 단계별로 기록: 발화율은 각성에서 최대, 서파 수면에서 낮고, 렘수면에서 거의 0이다 & \cite{AstonJonesBloom1981,McGintyHarper1976} & 측정됨 \\
4 & 렘수면에서 근육은 \nt{GABA}와 \nt{GLYC}를 통한 억제로 꺼진다 & 쥐에서 씹기 근육의 \nt{ACET} 뉴런 활동을 재고 차단제를 그 뉴런에 직접 주입: 두 유형의 \nt{GABA} 수용체와 \nt{GLYC} 수용체를 모두 차단해야만 마비가 풀린다 & \cite{BrooksPeever2012} & 측정됨 \\
5 & 수면 압력 $S$는 문턱값이 두 개인 커패시터다~\eqref{eq:sleeppressure}, $\tau_r=18.2$~h, $\tau_d=4.2$~h & 2과정 모델. 시상수는 EEG 서파 파워가 각성에 따라 커지고 수면 중에 줄어드는 방식에서, 문턱값의 모양은 자발적으로 깨는 시각에서 얻었다 & \cite{Borbely1982,Daan1984} & 이론 \\
6 & 기계는 스스로 $16.1$~h 각성과 $8.0$~h 수면에 이른다(그림~\ref{fig:twoprocess}) & 문턱값이 오르내리는 \eqref{eq:sleeppressure}에 따른 계산, 스크립트 \texttt{models/sleep\_models.py}: $16.1$~h 각성과 $7.9$--$8.0$~h 수면 & --- & 모델 \\
7 & $40$~h 동안 자지 않은 뒤 $S=0.90$이지만 수면은 $8.8$~h뿐이다. 빚은 길이가 아니라 깊이로 갚는다 & 모델: \texttt{models/sleep\_models.py}가 $S=0.90$과 $8.8$~h를 준다. 실험: $40.5$~h 동안 자지 않은 여덟 명에서 첫 회복 밤의 EEG 서파 파워가 평소보다 유의하게 높다 & \cite{Borbely1981} & 모델 \\
8 & 물리적으로 $S$는 아데노신이다 & 고양이 미세투석: 각성을 조절하는 영역 가운데 하나에서 세포 밖 아데노신이 깨어 있는 동안 오르고, 수면 박탈로 더 오르며, 자는 동안 떨어진다. $S$가 아데노신뿐이라는 것은 입증되지 않았다 & \cite{PorkkaHeiskanen1997,Dunwiddie2001} & 가설 \\
9 & 서파 수면에서 피질은 시소처럼 오간다: 수백 ms의 활동과 침묵, 1초에 한 번보다 드물게($<1$~Hz, 마취 상태에서 대개 $0.3$--$0.4$~Hz) & 마취한 고양이의 세포 안 기록: $0.8$--$1.5$~s의 탈분극과 긴 과분극, 리듬 $<1$~Hz, 대개 $0.3$--$0.4$~Hz. 같은 활동-침묵 교대가 자연 수면에서도 보인다(1행) & \cite{Steriade1993,Vyazovskiy2009} & 측정됨 \\
10 & 시소는 되먹임과 피로가 있는 무리다~\eqref{eq:updown}. $b=5$에서 주기 $1.1$~s(모델의 값), 활동은 시간의 약 $35\%$; $b=1$에서 고른 활동(그림~\ref{fig:updown}) & 계산, 스크립트 \texttt{models/sleep\_models.py}: 주기 $1.11$~s, 활동 시간 비율 $0.35$(정수 개 주기에 대한 평균); $b=1$에서 활동 비율 $1.0$. 모델의 주기는 측정값보다 짧다(9행) & --- & 모델 \\
11 & \nt{ACET}은 시소를 끄고 피질을 고른 작동으로 옮긴다 & 고양이에서 \nt{ACET} 뉴런 핵을 자극하자 피질 뉴런의 $79\%$에서 느린 리듬이 차단되고 고른 작동으로 바뀌었다. 주로 긴 과분극이 사라졌기 때문이다. 아세틸콜린은 피로를 만드는 느린 칼륨 전류를 억누른다 & \cite{SteriadeAmzica1993,McCormick1992} & 측정됨 \\
12 & $7$--$15$~Hz 리듬 폭발은 피질과 중계 핵 사이의 \nt{GLUT}--\nt{GABA} 루프가 만든다 & 페럿의 시각 중계 핵 절편에서, 폭발은 중계 세포가 스스로 흥분시키는 이웃 \nt{GABA} 핵의 억제에 대해 내는 반동 버스트로 생긴다 & \cite{vonKrosigk1993,SteriadeMcCormickSejnowski1993} & 측정됨 \\
13 & 하루의 재생은 빨리 감기로 진행된다: 해마에서 약 $20$배, 피질에서 $6$--$7$배 빠르다 & 서파 수면에서 해마 뉴런의 반복 순서가 시간적으로 압축되어 나타난다. 트랙을 달린 뒤 장소 뉴런의 순서가 같은 순서로, 약 $20$배 압축된 $\sim100$~ms의 짧은 폭발로 반복되었다. 피질에서는 $6$--$7$배 압축 & \cite{Nadasdy1999,LeeWilson2002,Euston2007} & 측정됨 \\
14 & 재생과 피질의 맞물림을 강화하면 기억이 좋아진다(인과 관계) & 쥐에서 학습 뒤 재생에 맞춘 자극으로 재생과 피질의 서파 및 리듬 폭발의 연결을 강화: 다음 날 기억이 높았고, 대조군은 우연 수준이었다 & \cite{Maingret2016} & 측정됨 \\
15 & 재생의 목적은 느린 기억의 섞어 배우기다 & 두 학습 시스템 이론과 인공 네트워크 계산: 순차 학습은 옛것을 망가뜨리고, 섞어 배우기는 그렇지 않다. 재생이 바로 이를 위해 필요하다는 뇌 직접 실험은 없다 & \cite{McClelland1995} & 가설 \\
16 & 잔 뒤 시냅스는 크기에 비례해 줄어들고, 큰 시냅스는 거의 바뀌지 않는다 & 생쥐 피질 시냅스 $6920$개의 3차원 전자 현미경: 잔 뒤 접촉 면적이 깨어 있은 뒤보다 $\sim18\%$ 작고, 감소는 크기에 비례하며, 크고 안정된 시냅스는 바뀌지 않는다 & \cite{deVivo2017} & 측정됨 \\
17 & 수면의 주된 과제는 가중치 재정규화다 & 시냅스 항상성 가설. 1행과 16행은 이와 부합하지만 이것이 주된 기능임을 증명하지는 않는다 & \cite{TononiCirelli2014} & 가설 \\
18 & 렘수면은 세계 모델의 벤치 시험이다 & 모드(표~\ref{tab:sleepmodes})는 측정되었다. 네트워크가 세계 모델을 돌린다는 것은 이론적 추측이며 실험은 없다 & \cite{Hobson2009} & 가설 \\
19 & 수면 중 뇌 세척: 열린 문제 & 한 실험: 잠자는 동안 세포 사이 공간이 $60\%$ 넓고 청소가 빠르다. 측정 방법이 다른 다른 실험: 잠자는 동안과 마취 상태에서 청소가 눈에 띄게 느리다. 결과가 서로 모순된다 & \cite{Xie2013,Miao2024} & 가설 \\
20 & 국소 수면: 깨어 있는 동물의 피질 일부가 잠깐 침묵으로 빠진다 & 오래 깨어 있은 쥐에서 피질 한 영역의 뉴런이 서파 수면에서처럼 잠깐 조용해지고, 국소 EEG에 서파가 나타난다. 그런 순간 쥐는 과제에서 더 자주 실수한다 & \cite{Vyazovskiy2011} & 측정됨 \\
\bottomrule
\end{longtable}}
